% Section 6
\section{Inference}\label{sec:inference}

\subsection{Sampling, estimators and regimes}

\paragraph{Where the randomness comes from.} \citet{aaiw2023} distinguish
sampling-based uncertainty (the units are a sample from a larger population) from
design-based uncertainty (some assignment is random). EI has neither: every unit
of the electorate is observed, and nothing is assigned. In a purely design-based
view nothing is random, $B$ is a fixed unknown, and the only valid bound is
Duncan--Davis. LCE, however, is a statement about a conditional expectation,
which has content only if unit rates are random given the unit's
characteristics. The uncertainty here is therefore \emph{model-based, conditional
on the design}, as in the fixed-regressor framework of \citet{aiz2014}:
\begin{itemize}
\item[(I1)] Given the design and the district effects, the unit rates $\beta_i$ are
independent with $\E[\beta_i\mid\text{design}]=m(V_i)$ (Assumption A6).
\item[(I2)] Inference is conditional on $\{(N_i,X_i,D_i)\}$, so the vertices of
the sample hulls are fixed.
\end{itemize}
District shocks are fixed features of this electorate. Under M1 and M2 they are
district effects, part of the mean. Under M0 a district shock is not a variance
problem but a misspecification: conditional on the design, M0's mean model holds
for the observed units only if the district shifts are zero
(Proposition~\ref{prop:tests}(c)), and clustered standard errors would not
repair it. No clustered variance is used. Independence given the district
effects is a substantive assumption, to be argued in each application: within-district
shocks can correlate residuals with no district mean shift, which violates (I1)
without violating the mean model. Residual correlation within districts and
between neighbours is a diagnostic, not a test that can confirm (I1). Under
dependence, the variance and residual-allowance statements below are not
covered.

\paragraph{Estimators.} Fit each of the $C-1$ free columns by OLS on the
homogeneous quadratic monomials $X_rX_s$ ($r\le s$). For M1 add the district shifts
$\one\{D=d\}X_r$ ($d\ge2$), which imposes the overidentifying restrictions of
Lemma~\ref{lem:D1}(b); for M2 fit one regression per district. Column $C$ follows
by adding up. Write $\hat\beta$ for the stacked coefficients, $\beta_0$ for their
true value and $\hat\Sigma$ for an HC3 covariance estimate. Each LP constraint is
$s_j(\beta)+a_j^\top k\ge0$, with $s_j$ linear in $\beta$ and $a_j$ fixed by the
design (Proposition~\ref{prop:lp}, Theorem~\ref{thm:ladder}(b)). The target is
$\Bt(k)=\varphi_0(\beta)+\omega^\top k$, with $\varphi_0$ linear in $\beta$ and
$\omega$ computed from the sample weights $w_{ir}$, so population moments are
never estimated. The \emph{plug-in bound} is
$\varphi_0(\hat\beta)+[\min,\max]\{\omega^\top k:\ s(\hat\beta)+Ak\ge0\}$. It is an
estimate of the identified set, not a confidence interval.

\begin{table}[t]
\centering\small
\begin{tabularx}{\linewidth}{@{}L{0.12\linewidth}L{0.15\linewidth}Y L{0.25\linewidth}@{}}
\toprule
Model & Regime & Result & Status \\
\midrule
any rung & any, any $n$ & Theorem~\ref{thm:finite} (finite-sample interval and emptiness test) & theorem; uninformative in width so far \\
M0 & any & Propositions~\ref{prop:consistency}, \ref{prop:triangular}, Theorems~\ref{thm:limit}, \ref{thm:ci} & theorem (asymptotic) \\
M1, M2 & R1 ($D^\star$ fixed) & the same & theorem (asymptotic) \\
M1, M2 & R3 & the same, with a high-dimensional Gaussian approximation~\citep{cck2017} under a rate condition & conjecture \\
M1 & R2 & coverage of the CLR interval & simulation finding \\
any & any & the CLR interval with simulated candidates & Conjecture~\ref{conj:clr} \\
\bottomrule
\end{tabularx}
\caption{Which inference result holds in which regime.}\label{tab:regimes}
\end{table}

Under M0 the number of constraints grows with the number of hull vertices, but
the regression has fixed dimension. Every constraint term is a linear functional
of one asymptotically normal vector, the events $\{\max_j(\cdot)\le c\}$ are convex
sets, and the multivariate central limit theorem holds uniformly over convex sets
\citep[Theorem~17.5]{bhattacharyarao1976}. So the growth in constraints does not
matter. Under M1 and M2 in R2 or R3, $\dim\beta$ grows with $D^\star$ and this
argument no longer applies.

\subsection{The structure of the bound and its plug-in estimator}

\begin{lemma}[An endpoint is a minimum of linear functions]\label{lem:minlin}\status{proved, fixed constraint system}
If the LP is feasible with bounded value at $\beta$, then
\[
\varphi_U(\beta)=\varphi_0(\beta)+\min_{\lambda\in\Lambda}\lambda^\top s(\beta),\qquad
\Lambda=\vertx\{\lambda\ge0:A^\top\lambda=-\omega\},
\]
a finite set. So $\varphi_U$ is concave and piecewise linear in $\beta$; likewise
$\varphi_L(\beta)=\varphi_0(\beta)-\min_{\lambda\in\Lambda'}\lambda^\top s(\beta)$ with
$\Lambda'=\vertx\{\lambda\ge0:A^\top\lambda=\omega\}$.
\end{lemma}
\begin{reading}
Each end of the bound is the tightest of finitely many candidate limits, each
linear in the regression coefficients; write $\theta_\lambda(\beta)=g_\lambda^\top\beta$
for the candidate with gradient $g_\lambda$. JKST's case has four candidates and
assumes one wins clearly. In R$\times$C with districts there are many candidates,
and near-ties are the norm. The lemma concerns a fixed constraint matrix and
objective, which conditioning on the design provides.
\end{reading}

\begin{proposition}[Consistency and rate]\label{prop:consistency}\status{proved}
Assume (I1)--(I2), A1--A5 for the rung used, $\sqrt n(\hat\beta-\beta_0)=O_p(1)$,
and Slater's condition: some $\tilde k$ and $\varepsilon>0$ satisfy
$s_j(\beta_0)+a_j^\top\tilde k\ge\varepsilon$ for every constraint defined by the
population hulls. Then (a) $\hat\varphi_U-\varphi_U(\beta_0;\hat A)=O_p(n^{-1/2})$,
and likewise for $\varphi_L$; (b) $\varphi_U(\beta_0;\hat A)\to\varphi_U(\beta_0;A)$
when $\hat A\to A$ in Hausdorff distance along a nested sequence, or along any
sequence whose feasible sets eventually lie in a common bounded set; (c) since
$\hat A_d\subseteq A_d$, $\varphi_U(\beta_0;\hat A)\ge\varphi_U(\beta_0;A)$ and
$\varphi_L(\beta_0;\hat A)\le\varphi_L(\beta_0;A)$.
\end{proposition}
\begin{reading}
The plug-in bound converges at the usual rate to the bound for the sample hull,
which is weakly wider than the population bound; estimating the hull only ever
widens it. The hull error is one-sided and tends to zero, but no rate is shown, so
the limit theory below concerns the sample-hull endpoint. Coverage of the true
parameter never depends on the hull, because the true parameter satisfies the
constraints at every observed composition. Slater's condition fails when every
feasible $k$ makes some implied rate exactly 0 at some vertex, a structural zero,
which is not exotic in EI; the rate result is then open, but the coverage results
below do not use Slater. Bounds do not shrink to a point: their width converges to
the population width.
\end{reading}

\begin{theorem}[Limit of the plug-in endpoint, fixed constraint system]\label{thm:limit}\status{cited structure}
Under the assumptions of Proposition~\ref{prop:consistency} and
$\sqrt n(\hat\beta-\beta_0)\Rightarrow\mathcal Z\sim N(0,\Sigma)$, suppose the
constraint system is fixed along the sequence, so that $A$, $\omega$, $\Lambda$ and
the gradients $g_\lambda$ do not depend on $n$. Let $\Lambda^*$ be the set of optimal
dual vertices at $\beta_0$. Then
\[
\sqrt n\big(\hat\varphi_U-\varphi_U(\beta_0)\big)\Rightarrow\min_{\lambda\in\Lambda^*}g_\lambda^\top\mathcal Z .
\]
(i) If $\Lambda^*$ is a singleton, the limit is normal. (ii) If two active
gradients differ in a direction in which the estimator varies,
$(g_\lambda-g_{\lambda'})^\top\Sigma(g_\lambda-g_{\lambda'})>0$, the limit has
strictly negative mean: the plug-in upper endpoint is biased downward, the lower
endpoint upward, and the plug-in bound is too narrow. If all active gradients
agree on the support of $\mathcal Z$, the limit is normal. (iii) In case (ii) the
standard bootstrap is inconsistent.
\end{theorem}
\begin{reading}
When several candidate limits are nearly tied, the plug-in end picks the one that
looks tightest by chance, so the estimated bound is biased inward. This is why the
plug-in is not a confidence interval. The proof applies the delta method for
directionally differentiable maps \citep{shapiro1991,fangsantos2019} and the
bootstrap characterisation of \citet{fangsantos2019}. With two tied candidates and
unit noise, the mean plug-in error is $-1/\sqrt\pi\approx-0.564$.
\end{reading}

\begin{remark}[Use $C-1$ columns]\label{rem:Cminus1}
The column shares sum to one, so the coefficients of the $C$ column regressions
sum to a known constant and the joint covariance of all $C$ columns is singular.
Stating $\hat\beta$, $\Sigma$ and the gradients for the first $C-1$ columns, and
recovering column $C$ by adding up, lets $\Sigma$ be non-singular. When it is
positive definite, the condition in (ii) reduces to $g_\lambda\ne g_{\lambda'}$.
\end{remark}

\begin{proposition}[Changing constraint systems]\label{prop:triangular}\status{proved}
Let the design, and hence the constraint system, change with $n$, with dual vertex
sets $\Lambda_n$, $|\Lambda_n|\le\bar L$, bounded gradients $g_{\lambda,n}$, and
localised slacks $c_{\lambda,n}=\sqrt n(\theta_{\lambda,n}(\beta_0)-\varphi_{U,n}(\beta_0))\ge0$.
Suppose that, after relabelling, $g_{\lambda,n}\to g_\lambda$ and
$c_{\lambda,n}\to c_\lambda\in[0,\infty]$. Then
\[
\sqrt n\big(\hat\varphi_{U,n}-\varphi_{U,n}(\beta_0)\big)\Rightarrow\min_{\lambda:\,c_\lambda<\infty}\big(g_\lambda^\top\mathcal Z+c_\lambda\big).
\]
The limit of Theorem~\ref{thm:limit} obtains when every $c_\lambda\in\{0,\infty\}$.
\end{proposition}
\begin{reading}
As the sample hull grows, candidates can move in and out of contention. A
candidate that is inactive only by a gap of order $n^{-1/2}$ still shapes the
limit. For example, $f_n(\theta)=\min\{\theta,-\theta+n^{-1/2}\}$ at $\theta_0=0$ has a
uniquely active first piece for every $n$, yet
$\sqrt n(f_n(\hat\theta)-f_n(\theta_0))\Rightarrow\min\{Z,1-Z\}$, which is not normal.
The result describes the sample-hull endpoint and says nothing about the
population-hull error. Uniform limit theory without convergent localised slacks is
open.
\end{reading}
\begin{sketch}
By Lemma~\ref{lem:minlin} the scaled error equals, exactly,
$\min_\lambda\big(g_{\lambda,n}^\top\sqrt n(\hat\beta-\beta_0)+c_{\lambda,n}\big)$.
Candidates with $c_{\lambda,n}\to\infty$ drop out with probability tending to one,
and the continuous mapping theorem gives the limit.
\end{sketch}

\subsection{Confidence intervals}\label{sec:ci}

Fix levels $\alpha_U,\alpha_L,\alpha_E\in(0,1)$. Let $\hat\sigma_\lambda$ be the
standard error of $\hat\theta_\lambda=\theta_\lambda(\hat\beta)$ computed from
$\hat\Sigma$, $\hat\sigma_s$ the vector of standard errors of the constraint terms
$s_j(\hat\beta)$, and $\hat\sigma_{\varphi_0}$ that of $\varphi_0(\hat\beta)$. For a cell
$(r,c)$ write $w_i=w_{ir}$, $S_1=\tfrac12(\sum_iw_i^2)^{1/2}$ (Popoviciu's bound
on the standard deviation of the residual term) and $z_{\alpha_E}=\Phi^{-1}(1-\alpha_E/2)$,
where $\Phi$ is the standard normal distribution function.

\begin{theorem}[Confidence intervals]\label{thm:ci}
\begin{enumerate}[label=(\alph*)]
\item \status{proved} \emph{Conservative interval.} Let $\hat c$ be the
$(1-\alpha_U-\alpha_L)$-quantile of the maximum of the standardised Gaussian errors of
all constraint terms (one-sided) and of the target constant (two-sided). Put
$\hat\Theta^+=\{k:s(\hat\beta)+\hat c\hat\sigma_s+Ak\ge0\}$ and
\[
\CI=\Big[\varphi_0(\hat\beta)-\hat c\hat\sigma_{\varphi_0}+\min_{\hat\Theta^+}\omega^\top k-z_{\alpha_E}S_1,\ \
\varphi_0(\hat\beta)+\hat c\hat\sigma_{\varphi_0}+\max_{\hat\Theta^+}\omega^\top k+z_{\alpha_E}S_1\Big].
\]
If $\max_iw_i^2/\sum_iw_i^2\to0$, then
$\liminf_nP(B_{rc}\in\CI)\ge1-\alpha_U-\alpha_L-\alpha_E$. Without the $S_1$ terms
the interval covers $\Bt_{rc}$ with probability at least $1-\alpha_U-\alpha_L$.
\item[(b1)] \status{cited} \emph{CLR with the full candidate set.} Let $\Lambda$ be
the full set of dual vertices of Lemma~\ref{lem:minlin}, finite for a fixed
design, and assume Condition~P of \citet{clr2013}: $\Sigma$ positive definite
(using $C-1$ columns and a full-rank design) and $g_\lambda\ne0$ for every
$\lambda$. For the upper end, select
$\hat V=\{\lambda:\hat\theta_\lambda\le\min_\mu(\hat\theta_\mu+\hat c_p\hat\sigma_\mu)+2\hat c_p\hat\sigma_\lambda\}$,
where $\hat c_p$ is the $1-0.1/\log n$ quantile of $\max_\lambda Z_\lambda$ and $Z$ is
the standardised Gaussian vector with correlations from $\hat\Sigma$; set
$\hat U=\min_{\lambda\in\hat V}(\hat\theta_\lambda+\hat c\hat\sigma_\lambda)$ with $\hat c$
the $(1-\alpha_U)$-quantile of $\max_{\lambda\in\hat V}Z_\lambda$; treat the lower end
alike; and add $z_{\alpha_E}S_1$ on each side. The coverage bound of (a) holds
\citep[Theorem~4]{clr2013}.
\item[(b2)] \status{proved} \emph{Any candidate set with a simultaneous critical
value.} Let $\hat\Lambda$ be any data-dependent set of dual-feasible points, and
$\hat U=\min_{\lambda\in\hat\Lambda}(\hat\theta_\lambda+\hat c\hat\sigma_\lambda)$, with
$\hat U=+\infty$ if $\hat\Lambda$ is empty. Then
$\liminf_nP(\varphi_U(\beta_0)\le\hat U)\ge1-\alpha_U$ in either of two cases:
(i) $\hat c$ is the $(1-\alpha_U)$-quantile of $\max_{\lambda\in\bar\Lambda}Z_\lambda$
for a fixed finite set $\bar\Lambda\supseteq\hat\Lambda$, with $\sigma_\lambda>0$ on
$\bar\Lambda$; (ii) \emph{projection}: $\hat c=(\chi^2_{q,1-\alpha})^{1/2}$ and
$\hat\sigma_\lambda=(g_\lambda^\top\hat\Sigma g_\lambda/n)^{1/2}$, where either
$\Sigma$ is positive definite and $q=\dim\beta$, or there is an exact reduced
full-rank parameterisation $\beta=L\eta+b$, $\hat\beta=L\hat\eta+b$, with $q=\dim\eta$,
$\Sigma_\eta$ positive definite and $\hat\Sigma=L\hat\Sigma_\eta L^\top$. In case
(ii) one event covers both ends, for any candidate sets, at total level $\alpha$.
\end{enumerate}
\end{theorem}
\begin{reading}
The intervals widen each candidate limit by its own uncertainty before taking the
tightest, so noisy candidates (for example, districts with few units) rarely set
the limit. Three versions carry proofs. The conservative interval widens all
constraints at once. The projection interval uses one ellipsoid for the whole
coefficient vector; it is valid for any collection of candidates, however chosen,
but it is wide when there are many coefficients, because its critical value grows
like $\sqrt q$. The fixed-superset version is sharper but needs a fixed superset.
For the actual rate $B$, an extra allowance covers each unit's local luck. JKST's
interval $\CI_x$ is version (b) with a single candidate per side and no selection
step, with the two error sources combined by adding standard deviations.

An exact reduction means documented identities, such as the adding-up identity
that removes column $C$, or coordinates fixed by the design. A rank read off
$\hat\Sigma$ with a numerical cutoff is not one. A singular limit does not control
errors of smaller order in its null directions: with
$\hat\beta-\beta_0=(Z/\sqrt n,-1/n)$, the second coordinate has standard error zero
and the bound $\hat\beta_2+0$ misses $\beta_{0,2}$ in every sample. So the projection
interval is stated with $q=\dim\beta$ unless an exact reduction is documented.
Dropping candidates that load on numerically null directions is then only a
safeguard: it can only widen the interval.
\end{reading}
\begin{sketch}
(a) On one simultaneous event, every true constraint term is at most its inflated
estimate and the target constant is within its allowance, so the true $k_0$ lies in
$\hat\Theta^+$ and $\Bt$ lies in the inflated LP range. The residual term $E$ has
conditional variance at most $S_1^2$ by Popoviciu's inequality; along
subsequences, Chebyshev's inequality handles the case where its actual standard
deviation is negligible relative to $S_1$, and the Lindeberg central limit theorem
the rest. A union bound finishes. (b2) By weak duality every dual-feasible
candidate is an upper bound for the true endpoint, so it suffices that all
candidates are simultaneously covered: by the fixed-superset maximum in (i), and by
Cauchy--Schwarz in the metric of $\hat\Sigma_\eta$ in (ii).
\end{sketch}

\begin{conjecture}[CLR with simulated candidates; the package default]\label{conj:clr}
Collect $\hat\Lambda$ as the optimal dual vertices of the LP at $\hat\beta$ and at
$n_{\rm draw}$ Gaussian draws around it, and compute both $\hat c_p$ and $\hat c$
over $\hat\Lambda$ in the rule of Theorem~\ref{thm:ci}(b1). Under the conditions
of (b1), the resulting interval has asymptotic coverage at least $1-\alpha_U$ for
the upper end (and likewise for the lower end) as $n_{\rm draw}\to\infty$.
\end{conjecture}

The procedure is an \emph{algorithm}: its validity is not proved. What is proved
is a reduction.

\begin{lemma}[Reduction to the full-candidate rule]\label{lem:reduction}\status{proved}
Suppose $\hat c_p$ is computed over the full set $\Lambda$, and $\hat\Lambda$
contains every vertex selected by the rule of Theorem~\ref{thm:ci}(b1) over
$\Lambda$ and a minimiser of $\hat\theta_\mu+\hat c_p\hat\sigma_\mu$ over $\Lambda$.
Then the threshold, the selected set, $\hat c$ and $\hat U$ coincide with those of
(b1), so its coverage carries over.
\end{lemma}
\begin{reading}
Three steps remain open: recovery of vertices that are optimal only for rare
Gaussian directions (and degenerate optimal faces); the fact that computing
$\hat c_p$ over a subset of $\Lambda$ can only lower it, an error in the
anti-conservative direction; and the growth of $|\Lambda|$ with the number of
districts, outside CLR's fixed-index setting. Until these are settled, the CLR
interval as implemented is an algorithm with simulation support
(Section~\ref{sec:evidence}), and the intervals of Theorem~\ref{thm:ci}(a) and
(b2) carry the formal guarantee.
\end{reading}

\paragraph{A variance-adaptive residual allowance.} Popoviciu's bound uses the
worst-case variance $\tfrac14$ for every unit. The Bhatia--Davis inequality
\citep{bhatiadavis2000} gives $\Var(e_{irc}\mid V_i)\le m_i(1-m_i)$ with
$m_i=m_{rc}(V_i)$. Put $\pi_i=w_i^2/\sum_jw_j^2$ and $\bar m=\sum_i\pi_im_i$, the
$\pi$-weighted model-implied rate, which is itself a linear target. By concavity of
$m(1-m)$,
\[
s_n^2:=\Var(E_{rc}\mid\text{design})\le\sum_iw_i^2m_i(1-m_i)\le\Big(\sum_iw_i^2\Big)\bar m(1-\bar m)=:\bar s^2 .
\]
Let $\CI_v$ be a confidence interval for $\bar m$ at level $1-\alpha_v$, and
$\hat S_{BD}=\big(\sum_iw_i^2\max_{m\in\CI_v\cap[0,1]}m(1-m)\big)^{1/2}$.

\begin{proposition}[Variance-adaptive allowance]\label{prop:bd}\status{proved}
Suppose $\CI_v$ is an interval of Theorem~\ref{thm:ci}(a) or a projection interval
of Theorem~\ref{thm:ci}(b2)(ii), computed on the rung's own constraint system
(without added sign rows), and $\max_iw_i^2/\bar s^2\to0$. Replacing $z_{\alpha_E}S_1$
by $z_{\alpha_E}\hat S_{BD}$ in the conservative interval, or in the projection
interval for $B$, gives
$\liminf_nP(B_{rc}\in\CI)\ge1-\alpha_{\rm est}-\alpha_E-\alpha_v$, where
$\alpha_{\rm est}$ is the level spent on the estimation part.
\end{proposition}
\begin{reading}
The allowance for local luck can be made smaller when the plausible rates are far
from one half, provided the input interval for $\bar m$ comes from a proved
method. The condition on the weights holds, for instance, when $\bar m$ stays
bounded away from 0 and 1. If $\CI_v$ fails, the allowance is not conservative:
with true $\bar m=.5$, a wrong interval inside $[0,.1]$ gives a variance bound of at
most $.09$ against the true $.25$. If $\CI_v$ is instead the CLR interval, the
same bound holds conditional on $\CI_v$ covering $\bar m$, and the result inherits
the status of Conjecture~\ref{conj:clr}. An optional, stronger assumption (the
unit-level deviations of different groups are not negatively correlated) would
make the residual variance estimable from the outcome shares and remove most of
the allowance; it is not in the default.
\end{reading}
\begin{sketch}
On the event that $\CI_v$ covers $\bar m$, $\hat S_{BD}\ge\bar s\ge s_n$. So a
miss of the residual allowance requires either that $\CI_v$ fails or that
$|E|>z_{\alpha_E}\bar s$, whose limiting probability is at most $\alpha_E$ by the
Chebyshev--Lindeberg argument with $\bar s$ in place of $S_1$. A union bound with
the estimation event finishes.
\end{sketch}

\paragraph{Boundary cases.} When the true $K$ is on the boundary of the identified
set, the side of the interval that touches it carries the full error probability,
and coverage is close to nominal. When the true $K$ is interior, coverage tends to
one: this is the distinction between covering the parameter and covering the
set~\citep{imbensmanski2004}. When the bound's width is near zero (L4 in
$2\times2$, L5), one-sided intervals per endpoint under-cover unless the critical
value adapts~\citep{stoye2009}. The two-sided split used here is uniformly valid
at some cost in width.

\subsection{Tests of the model}\label{sec:tests}

\begin{proposition}[Tests]\label{prop:tests}
(a) \status{proved} \emph{Emptiness test.} Rejecting the rung when $\hat\Theta^+$ is
empty has asymptotic level at most $\alpha_U+\alpha_L$.
(b) \status{proved for fixed dimension} \emph{Shared against district-specific
effects.} A Wald test of the M1 regression against M2 tests the
$(D^\star-1)(C-1)R(R-1)/2$ overidentifying restrictions of Lemma~\ref{lem:D1}(b);
its $\chi^2$ calibration holds for fixed dimension (regime R1).
(c) \status{proved} \emph{Pooling.} Conditional on the design, suppose each
district's observed compositions contain $R$ affinely independent points. Under
M1, M0's mean model holds for the observed units iff all district shifts vanish.
Pooling is therefore tested by a Wald test of $\zeta=0$ in the M1 regression, with
$(D^\star-1)(C-1)R$ restrictions.
\end{proposition}
\begin{reading}
An empty confidence set is a real rejection at the stated level. The curvature
test checks the testable part of shared effects. The pooling test says whether the
pooled model is wrong for this design; a rejection does not mean that its
intervals fail (Section~\ref{sec:evidence}), and the test has no power against
within-district dependence without mean shifts. Proposition~\ref{prop:D4}(b) gives a
weaker condition for pooling, but it averages over district labels and so
describes a superpopulation of re-drawn districts, which is not the target here;
conditional on the observed labels only $\zeta=0$ makes M0 correct.
\end{reading}
\begin{sketch}
(a) On the simultaneous event of Theorem~\ref{thm:ci}(a), the true $k_0$ lies in
$\hat\Theta^+$, so emptiness requires that event to fail. (c) If M0 held for the
observed units with coefficients $H'$, then in each district
$(G_c-\zeta_{c,d}\one^\top)X_i=0$ for all its units, with $G=H'-H$. The district's
$R$ affinely independent compositions lie in $\{\one^\top x=1\}$ and are therefore
linearly independent, so $G_c=\zeta_{c,d}\one^\top$ in every district. With
$\zeta_{c,1}=0$, $G_c=0$ and every shift vanishes.
\end{sketch}

\paragraph{Calibration by a restricted wild bootstrap.} In moderate dimension the
HC3 Wald statistic compared with $\chi^2$ over-rejects. The package calibrates
both tests by a restricted wild bootstrap: fit the null model (M0 for the pooling
test, M1 for the curvature test); draw $Y^*=\hat Y_0+\hat e_0v$ with Rademacher
$v_i$, one per unit, shared across the $C-1$ equations; refit the larger model and
recompute the same Wald statistic; and report
$p=(1+\#\{W^*\ge W\})/(B+1)$ with $B=399$ by default. The $\chi^2$ and F p-values are
reported as well. This calibration is supported by simulation
(Section~\ref{sec:evidence}); its validity as the number of restrictions grows is
not proved.

\subsection{A finite-sample interval valid in every regime}

Fix a rung and condition on the design. Let $Z$ be the rung's regression design,
with full column rank, and $\hat\beta_c=(Z^\top Z)^{-1}Z^\top T_{\cdot c}$ for
$c<C$. Every linear functional $f(\beta)=f_0+\sum_{c<C}b_c^\top\beta_c$ satisfies
$f(\hat\beta)=f_0+\sum_i\sum_{c<C}w_{ic}(f)T_{ic}$ with
$w_{ic}(f)=b_c^\top(Z^\top Z)^{-1}z_i$. Put $w_{iC}(f)=0$,
$\rho_i(f)=\max_cw_{ic}(f)-\min_cw_{ic}(f)$, and the Hoeffding radius
$h(f;\delta)=\big(\tfrac12\log(1/\delta)\sum_i\rho_i(f)^2\big)^{1/2}$. With $J$
constraint rows and $\delta=\alpha_s/(J+2)$, put $h_j=h(s_j;\delta)$,
$h_0=h(\varphi_0;\delta)$, $t_E=\big(\tfrac12\log(2/\alpha_E)\sum_iw_{ir}^2\big)^{1/2}$,
$\hat\Theta^{\rm fin}=\{k:s_j(\hat\beta)+h_j+a_j^\top k\ge0\ \forall j\}$, and
\[
\CI^{\rm fin}=\Big[\varphi_0(\hat\beta)-h_0+\min_{\hat\Theta^{\rm fin}}\omega^\top k-t_E,\ \
\varphi_0(\hat\beta)+h_0+\max_{\hat\Theta^{\rm fin}}\omega^\top k+t_E\Big].
\]

\begin{theorem}[Finite-sample interval and emptiness test]\label{thm:finite}\status{proved}
Assume (I1), that the rung's mean model holds and that $Z$ has full column rank.
Then for every $n$ and in every regime: (a)
$P(B_{rc}\in\CI^{\rm fin}\mid\text{design})\ge1-\alpha_s-\alpha_E$; (b) without $t_E$
the interval covers $\Bt_{rc}$ with probability at least $1-\alpha_s$; (c)
rejecting the rung when $\hat\Theta^{\rm fin}$ is empty is a test of exact level at
most $\alpha_s$; (d) intersecting with Duncan--Davis keeps (a).
\end{theorem}
\begin{reading}
This interval needs no Gaussian approximation, no fixed dimension, no Slater
condition and no regime, so it is the only formal guarantee for shared effects with
many small districts. In every design tried it is as wide as the logical range
$[0,1]$. Two costs combine: Hoeffding's inequality uses the worst-case variance
for each outcome share (the true variance was about 0.01, against the worst case
$\tfrac14$), and the
constraint-by-constraint union bound inflates each radius, which is then divided
by the distance from the vertex to a pure composition. An oracle Gaussian union
bound with the true variances still gave widths of 0.76--0.87, so most of the cost
is the union bound, which the conservative asymptotic interval shares. The theorem
settles validity in regime R2, but not usefulness; it is not in the package.
\end{reading}
\begin{sketch}
Every constraint term and the target constant are linear functionals of the OLS
fit, and each is a sum over units of independent terms whose range is known from
the design, because the outcome shares lie in the simplex. OLS is conditionally
unbiased under the mean model, so Hoeffding's inequality~\citep{hoeffding1963}
gives a finite-sample radius for each. A union bound over the $J+2$ one-sided
events puts the true $k_0$ in the widened set, and the residual term gets its own
Hoeffding radius.
\end{sketch}

\subsection{Many districts with few units}

\begin{proposition}[Dropping and selecting districts]\label{prop:select}\status{proved}
(a) For any subset $\mathcal D'$ of districts, the M1 outer set computed from the
constraints of $\mathcal D'$ contains the full M1 outer set. (b) If $\mathcal D'$ is
chosen from the design only, any procedure with (finite-sample or asymptotic)
coverage for a fixed set of districts keeps that coverage conditionally on the
selected $\mathcal D'$; in particular Theorem~\ref{thm:finite} holds. The asymptotic
results transfer only if their own regime conditions hold for the retained
districts.
\end{proposition}
\begin{reading}
With many small districts, ignoring some districts' constraints can only widen the
bound, so it never costs validity of the identified set. If the retained districts
are chosen from the design alone (for example, those with at least 30 units or
enough spread in composition), not from the outcomes, any interval that is valid
for a fixed set of districts stays valid. Selection does not by itself move the
problem into a regime where the asymptotic theory applies: retained districts with
few units are still in R2. Other practical options are M1($\delta$) with $\delta>0$,
which bounds how far one district can pull the shared parameter, coarser groups,
and targeting $\Bt$ to drop the residual term. A hierarchical model for the district
shifts would change the model and is not pursued.
\end{reading}

\subsection{Sensitivity to nonlinearity}\label{sec:sensitivity}

\begin{proposition}[Bounded nonlinearity]\label{prop:S1}\status{proved; an outer bound}
Replace (S2) by $\E[\beta\mid V]=M_H(X)+\nu(X)$ with $\nu$ in the undetectable class
of Proposition~\ref{prop:undetectable} ($x^\top\nu(x)=0$, $\nu(x)\one=0$) and
$\|\nu(x)\|_\infty\le\eta$. Then $Q$ is identified as before, and every value of
$\Bt_{rc}$ consistent with the data lies in
\[
\Big[\min_{\ThO(\eta)}\varphi-\eta,\ \max_{\ThO(\eta)}\varphi+\eta\Big],\qquad
\ThO(\eta)=\{k:s+Ak\ge-\eta\one\}.
\]
For a general target $\sum a_{rc}B_{rc}$, widen by $\eta\sum|a_{rc}|$ instead.
\end{proposition}
\begin{reading}
If rates may depart from straight lines, invisibly, by up to $\eta$, the bound
widens in a controlled way. The bound is valid, but it is not claimed to be the
sharp set of the relaxed model. The \emph{breakdown value} $\eta^*$ of a
conclusion, such as ``loyalty exceeds 0.30'', is the largest departure under which
it survives; the bounds are nested in $\eta$, so it can be found by bisection. In
the $3\times3$ illustration of Section~\ref{sec:rxc}, $\eta^*(\text{loyalty}>0.30)=0.060$:
the conclusion survives any undetectable nonlinearity of up to six percentage
points. Applying the confidence-interval version of the bound accounts for
sampling error. The dial $\delta$ of Proposition~\ref{prop:D6} works the same way
for heterogeneity across districts.
\end{reading}

\subsection{Inference with sign rows, unions and near-empty sets}\label{sec:R10}

The sign rows of Section~\ref{sec:msign} enter the same machinery. A row
$s_j(\beta)+a_j^\top k\ge0$ with $s_j$ linear in $\beta$ and $a_j$ fixed by the design
is \emph{linear in $\beta$}; every sign, ordering and local row has this form,
because the associations, the targets and $\ER^*$ are linear in the model means
with design weights, and $Q$ is linear in $\beta$. A \emph{parallel disjunction} is
a row $\max_{l\le L}s_l(\beta)+a^\top k\ge0$, whose pieces share the normal $a$. A
\emph{general union} is $C_+\cup C_-$, each branch a finite set of rows linear in
$\beta$. $\CI(\mathcal R)$ denotes an interval of Theorem~\ref{thm:ci} computed with the
base rows plus the rows $\mathcal R$. For version (a), $\CI=\emptyset$ iff
$\hat\Theta^+=\emptyset$; for version (b), iff the computed ends cross.

\begin{theorem}[Inference with added rows]\label{thm:R10}\status{proved for pooled M0, by interval version}
Assume the conditions of Theorem~\ref{thm:ci} for the augmented row system, the
pooled model M0, and that the restrictions hold at the truth.
\begin{enumerate}[label=(\alph*)]
\item \emph{Added rows.} For finitely many rows linear in $\beta$,
Theorem~\ref{thm:ci} and Proposition~\ref{prop:tests}(a) hold for the augmented
system, version by version (Table~\ref{tab:R10}). For an actual target that is a
contrast, such as $B_1-B_0$, the residual allowance must cover the contrast:
$\tfrac12\{\sum_i(w_{i1}+w_{i0})^2\}^{1/2}$ is a valid standard-deviation envelope,
and so is the sum of the per-cell envelopes.
\item \emph{Parallel disjunction} (conservative version only). Inflate each piece
separately and then take the maximum,
$\max_l\big(s_l(\hat\beta)+\hat c\hat\sigma_l\big)+a^\top k\ge0$, with all pieces in the
critical-value maximum. This covers at the nominal level. One-vs-rest and
$2\times2$ CR are of this form (the \emph{slab}), and no branch is selected from
$\widehat{\ER}$.
\item \emph{General union.} $\CI(C_+)\cup\CI(C_-)$, each branch at the full level
$\alpha$, covers the true target with asymptotic probability at least $1-\alpha$. No
Bonferroni split is needed. This is point coverage; it gives no simultaneous
coverage of both branch images.
\item \emph{Near-empty sets.} Under the model,
$\limsup_nP(\CI=\emptyset)\le\alpha_U+\alpha_L$, including point-identified cases,
and conditioning on non-emptiness can only raise coverage. If the restrictions are
false, no coverage statement holds.
\item \emph{Plug-in brackets are not valid.} Imposing the estimated bracket
$\hull\{\gamma_r\widehat{\ER},\widehat{\ER}\}$ without inflating it can under-cover,
and more data do not cure it: in a repeated three-point design,
$\limsup_nP(\tilde D_0\in\CI_{\rm naive})\le\Phi(\tfrac7{12}z_{.995})=0.9335<0.95$.
\end{enumerate}
\end{theorem}

\begin{table}[t]
\centering\small
\begin{tabularx}{\linewidth}{@{}L{0.32\linewidth}Y L{0.17\linewidth}@{}}
\toprule
Interval version & What Theorem~\ref{thm:R10} claims & Status \\
\midrule
(a) conservative & (a)--(d), including the slab (b) & proved \\
(b1) CLR over all dual vertices & (a), (c), (d), under CLR's conditions for the full augmented system & cited \\
(b2)(i) fixed superset & the branch-union and crossing arguments & proved \\
(b2)(ii) projection & both ends, any candidates; $\Sigma$ positive definite or an exact reduction & proved \\
CLR with simulated candidates (package default for CR) & the inclusion and crossing steps are deterministic; their probabilities are conjectural & algorithm (Conjecture~\ref{conj:clr}) \\
\bottomrule
\end{tabularx}
\caption{Status of Theorem~\ref{thm:R10} by interval version.}\label{tab:R10}
\end{table}

\begin{reading}
Sign beliefs enter the same interval machinery as the rung's own constraints.
Reinforcement uses a band whose two pieces are each widened by their own
uncertainty before the larger is taken; this is valid without knowing which branch
is true. The tempting shortcut, plugging the estimated neighbourhood and regression
numbers into the bracket, is not a confidence interval. Intersecting the fitted
bracket with a finished JKST interval is worse: its limiting coverage of a truth on
the bracket's edge is at most one half. When the belief is false, the interval stays
narrow and can be reproducibly wrong (Section~\ref{sec:evidence}); a narrow,
non-empty interval is no evidence that the restrictions are true. Conditioning on
other data events, or selecting a belief because it gives an attractive result, is
not covered.
\end{reading}
\begin{sketch}
(a) The added rows are linear in the regression coefficients, so they join the same
simultaneous Gaussian event, on which the true $k_0$ satisfies every inflated row.
(b) The piece that holds at the truth is dominated on the same event, so no branch
is ever chosen from the data. (c) The true branch is a fixed population fact; the
interval for that branch alone has the full level and lies inside the union.
(d) An empty interval requires the coverage event to fail. (e) In the three-point
design, with the basis $((1-x)^2,2x(1-x),x^2)$, the target constant's standard error
is exactly $\tfrac7{12}$ of the regression slope's, and the critical value is at
most $z_{.995}$ by the union bound over the eight base rows and the two target
tails.
\end{sketch}

The statements are pointwise in the data-generating process. Uniform validity near
point identification is open. Statements of the type of
Proposition~\ref{prop:consistency} and Theorem~\ref{thm:limit} need Slater's
condition, which rows that bind at the truth (such as the edge of the CR band) can
break; the conservative version does not.

\subsection{A practical recipe}

\begin{enumerate}
\item Fit the rung's regression on $C-1$ columns with HC3 covariance, and compute
the sample hulls.
\item Build the LP and report the plug-in bound, labelled as an estimate of the
identified set.
\item Report the CLR interval for $\Bt$ and for $B$, labelled as algorithmic, and
beside it the projection interval as the interval with an asymptotic guarantee;
for $B$, take $\bar m$ from the projection interval on the rung alone, or use
Popoviciu's bound. Intersect both with Duncan--Davis, which always contains $B$.
\item With many small districts, say that the CLR interval's coverage rests on
simulation. Theorem~\ref{thm:finite} gives a finite-sample guarantee, but so far
only at the width of the logical range.
\item Report the emptiness test and, for M1, the curvature and pooling tests with
wild-bootstrap p-values. A rejection of pooling says that M0 is wrong for this
design, not that its intervals fail.
\item Report sensitivity curves in $\eta$ and $\delta$, with the breakdown value of
the headline conclusion.
\item With few units per district, select districts by design, or prefer
M1($\delta$) or coarser groups, and say so.
\end{enumerate}
